\documentclass[12pt,a4paper]{article}
\usepackage[utf8]{inputenc}
\usepackage[T1]{fontenc}
\usepackage[spanish,es-noquoting]{babel}
\usepackage{geometry}
\usepackage{xcolor}
\usepackage{enumitem}

\geometry{a4paper, margin=1in}

\title{Evaluación de Examen Teórico-Práctico}
\author{Prof. César Rodríguez \\ \small Curso: Seguridad Informática}
\date{\today}

\begin{document}

\maketitle

\section*{Datos del Estudiante}
\textbf{Nombre:} Daniel Albadawi \\
\textbf{Grupo:} 2 (OSINT y Huella Digital) \\
\textbf{Calificación Final:} \textbf{\textcolor{blue}{10/10}}

\section*{Evaluación Detallada}

\begin{itemize}[label=$\bullet$]
    \item \textbf{Pregunta 1 (Significado de OSINT):} \textbf{2/2 puntos}. La definición es exacta y captura la esencia principal del OSINT, que es la legalidad y el uso de fuentes públicas en contraste con el espionaje o la inteligencia clasificada.
    
    \item \textbf{Pregunta 2 (Datos críticos en WHOIS):} \textbf{2/2 puntos}. Identificó correctamente los tres bloques de información más críticos que un auditor busca en un registro WHOIS durante la fase de reconocimiento de un dominio (servidores de nombres, contactos y fechas).
    
    \item \textbf{Pregunta 3 (Google Dork \texttt{site:dominio.com -www}):} \textbf{2/2 puntos}. Demuestra entendimiento de cómo funciona el operador lógico de exclusión (\texttt{-}) y comprende el valor táctico en una auditoría (descubrir subdominios ocultos o entornos vulnerables).
    
    \item \textbf{Pregunta 4 (Búsqueda de archivos \texttt{.ini} o \texttt{.env}):} \textbf{2/2 puntos}. La sintaxis de los operadores de Google es perfecta. Utiliza correctamente el operador \texttt{OR} y restringe la búsqueda al dominio objetivo. Acertada la anotación sobre \texttt{filetype:}.
    
    \item \textbf{Pregunta 5 (Error HTTP 429 en Python):} \textbf{2/2 puntos}. Excelente respuesta. Identifica el \textit{Rate Limiting} (demasiadas peticiones) y ofrece soluciones programáticas estándar y efectivas en la industria para el \textit{web scraping} y la automatización de OSINT (\texttt{time.sleep()}, \textit{exponential backoff}, rotación de proxies).
\end{itemize}

\vspace{1cm}

\section*{Comentarios Finales}
\noindent El estudiante demuestra un dominio sólido y sobresaliente de los conceptos de OSINT y el uso de Google Dorks. Sus respuestas son claras, concisas y denotan una comprensión tanto teórica como práctica, lo cual es vital para el desarrollo del módulo de OSINT dentro de la Suite de Auditoría.

\vspace{0.5cm}
\begin{center}
    \textit{¡Excelente trabajo!}
\end{center}

\end{document}